\subsection{The exact archimedean constant of the key rule}
\label{subsec:ftb-key-archimedean-constant}

We use the actual complex uniformization comparison in the following
form.  If \(x/\mathbb C\) is a key datum and \(u\) any of its
holomorphic line uniformizations with Appell--Humbert data
\((H,\chi)\), then \(E=\operatorname{Im}H\) on the period lattice
\(\Lambda\) is integral and induces an isomorphism
\(\Lambda\to\Lambda^\vee\).
Every algebraic tangent vector is the derivative of a unique
analytic vector and conversely; this comparison respects linear
combinations.
This is the precise fixed-fibre input
\srcid{unimodular_and_tangent_of_isAH}.
It concerns this actual line bundle and this actual \(u\);
it is not an assertion about an independently chosen lattice
or Hermitian form.

For clarity, this comparison transfers from the principal
uniformization supplied by the fixed-fibre realization to any
\(u\) as follows.
The analytic exponential maps have a unique complex-linear lift
\(F\) between their covers.
It is injective: if \(Fz=0\), then every \(tz\), \(t\in\mathbb C\),
belongs to the discrete period lattice; continuity makes \(z=0\).
It is therefore invertible in equal dimension and identifies
their lattices.
The quotient of the two local line frames glues to a nowhere-zero
entire function.
The Appell--Humbert triviality calculation applied to this quotient
identifies the two Hermitian forms after pullback by \(F\).
Transport the perfect alternating pairing through the resulting
lattice isomorphism.
The chain rule transports the tangent comparison.
These are exactly the two additional comparisons needed when
the given \(u\) is not the original chosen uniformization.
Their actual analytic realization is retained with the
fixed-fibre producer.

\begin{srlemma}[Canonical theta integral]
\label{lem:ftb-key-canonical-theta-integral}
Let \(\tau=\tau^t\), \(Y=\operatorname{Im}\tau>0\), and
\(a,b\in\mathbb R^g\).  For the source normalization
\[
 \theta[a,b](z,\tau)=
 \sum_{n\in\mathbb Z^g}
 e^{\,\pi i(n+a)^t\tau(n+a)+2\pi i(n+a)^t(z+b)}
\]
one has
\[
 \int_{[0,1]^{2g}}
 |\theta[a,b](x+\tau y,\tau)|^2e^{-2\pi y^tYy}\,dx\,dy
                           =(2^g\det Y)^{-1/2}.
 \tag{A1}\label{eq:ftb-key-theta-integral}
\]
Replacing closed cubes by half-open cubes does not change the
integral.
\end{srlemma}
\begin{proof}
The Gaussian decay in \(n\) makes the theta series normally
convergent on compact sets.
For fixed \(y\), its Fourier coefficients in \(x\), after
the harmless modulus-one factor \(e^{2\pi i a^tx}\), are
absolutely summable.
Parseval, or integration of the absolutely convergent double
series, therefore gives
\[
 \int_{[0,1]^g}|\theta[a,b](x+\tau y,\tau)|^2
                         e^{-2\pi y^tYy}\,dx
       =\sum_{n\in\mathbb Z^g}
                    e^{-2\pi(y+n+a)^tY(y+n+a)}.
\]
The terms are nonnegative, so Tonelli applies without a
prior integrability assertion.
Unfold the translated unit cubes to \(\mathbb R^g\),
then translate by \(a\).
Write \(Y=B^tB\).
The real Jacobian is \(|\det B|=(\det Y)^{1/2}\);
the remaining integral is the product of the one-variable
integrals \(\int_{\mathbb R}e^{-2\pi t^2}dt=2^{-1/2}\).
This proves \eqref{eq:ftb-key-theta-integral}, including its
finiteness.
Cube boundaries have measure zero.
\end{proof}

\begin{srproposition}[The exact key-rule norm]
\label{prop:ftb-key-exact-archnorm}
Let \(x/\mathbb C\) be an actual key datum of dimension \(g\).
Let \(u\), \(H,\chi\), \(\theta\), and \(W=(W_i)_i\) satisfy:
\(u\) is its holomorphic uniformization with actual
Appell--Humbert data \((H,\chi)\), respects its origin
rigidification, reads its theta section as \(\theta\), and
has derivative \(W_i\) equal to the \(i\)-th vector of the
ordered algebraic Lie frame.
For any additive fundamental domain \(\mathcal F\) for its
period lattice, set
\[
 n_\theta=
 \frac{\displaystyle\int_{\mathcal F}|\theta(z)|^2
                           e^{-\pi H(z,z)}\,dz}
      {\operatorname{vol}\mathcal F},
 \qquad
 n_\eta=
 \frac{\operatorname{vol}\mathcal F}
      {\pi^g|\det W|^2}.
\]
Under the fixed-fibre comparison stated above, the
exponent-one rule of
Theorem~\ref{thm:ftb-key-rule-exponent-one} satisfies
\[
                       |\kappa_{\mathbb C}(x)|\,
                                  n_\theta^4n_\eta^2=(2\pi)^{2g}.
 \tag{A2}\label{eq:ftb-key-archnorm}
\]
All factors refer to the displayed \(u,H,\theta,W\) and
\(\mathcal F\).
\end{srproposition}
\begin{proof}
The tangent comparison makes \(W\) invertible: a relation
among its rows would give the same relation among an
algebraic Lie basis, and conversely every analytic
vector is the derivative of a tangent vector.

Choose a symplectic integral basis \(e_i,f_i\) for \(-E\).
The full lattice rank is \(2g\), since it is discrete and
spans \(\mathbb C^g\) over \(\mathbb R\).
Let \(P\) have columns \(e_i\).
The Gram matrix \(G=P^tH\overline P\) is real symmetric,
because \(E(e_i,e_j)=0\), and is positive definite.
The \(e_i\) are real independent as part of the real
basis of a full lattice; hence positivity first gives
positive definiteness over \(\mathbb R\), then over
\(\mathbb C\).  Thus \(P\) is invertible.
Put \(\phi(z)=Pz\) and take the columns of \(\tau\) to
be \(P^{-1}f_i\).
The symplectic relations give
\[
 G\,\operatorname{Im}\tau=I,\qquad
 \tau=\tau^t,\qquad \operatorname{Im}\tau>0.
\]
Indeed the first equality is
\(\operatorname{Im}H(e_i,f_j)=-\delta_{ij}\);
the \(f_i,f_j\) isotropy relations then give the
symmetry of the real part of \(\tau\).
Consequently, writing \(Y=\operatorname{Im}\tau\),
\[
 \phi^{-1}\Lambda=\mathbb Z^g+\tau\mathbb Z^g,
 \qquad H(\phi z,\phi w)=z^tY^{-1}\overline w.
 \tag{A3}\label{eq:ftb-key-ah-coordinate-change}
\]

Symmetry of the actual line makes \(\chi(l)^2=1\).
Here is the relevant line calculation.
The inverse-translation automorphy laws first give
\(\chi(l)\chi(-l)=1\).
The line isomorphism under inversion gives a nowhere-zero
entire quotient with multiplier \(\chi(-l)/\chi(l)\).
These multipliers have modulus one, so the quotient's
absolute value is full-lattice periodic.
It is bounded on a compact real period parallelepiped,
hence the quotient is bounded everywhere and is constant
by Liouville.
Thus \(\chi(-l)=\chi(l)\), proving the assertion.

The actual section reading is entire and has the
Appell--Humbert automorphy law.
It is nonzero because the section is fibrewise nonzero
and the exponential covers all complex points.
Multiplying \(\theta(\phi z)\) by
\(e^{-\pi z^tY^{-1}z/2}\) changes that law to the classical
theta law with a sign semicharacter.
The semicharacter's values on the integral symplectic
basis give \(a,b\in\{0,\tfrac12\}^g\).
The unit-period Fourier expansion and the \(\tau\)-period
recurrence determine every coefficient from one
coefficient.  Namely the coefficient with index \(n+a\)
is that coefficient times
\(e^{\pi i(n+a)^t\tau(n+a)+2\pi i(n+a)^tb}\),
with the constant coefficient absorbed into the scalar.
Hence, for some \(c\ne0\),
\[
 \theta(\phi z)=
       c\,e^{\pi z^tY^{-1}z/2}\theta[a,b](z,\tau).
 \tag{A4}\label{eq:ftb-key-gaussian-gauge}
\]
The series is normally convergent, so coefficient comparison
and the recurrence are valid as identities of entire
functions.  Change the line frame by the indicated
Gaussian.  Its value at zero is \(1\), so the given origin
rigidification is preserved.

Define the rows of \(M\) by
\[
                         M_i=2\pi i\,\phi^{-1}(W_i).
\]
This is the actual tangent normalization after the
coordinate and Gaussian frame change.
It gives \(M\) invertible, and this transformed
uniformization is a normal form for the same key datum.
Therefore \(\kappa_{\mathbb C}(x)=\det(M)^4/c^8\).

The real period map
\((x,y)\mapsto\phi(x+\tau y)\) has determinant
\(|\det\phi|^2\det Y\): in real and imaginary coordinates
its matrix is the product of \(\phi\) with the block
matrix \(\left(\begin{smallmatrix}I&\operatorname{Re}\tau\\0&Y
\end{smallmatrix}\right)\).
Thus
\[
                    \operatorname{vol}\mathcal F
                                   =|\det\phi|^2\det Y.
 \tag{A5}\label{eq:ftb-key-period-volume}
\]
To use an arbitrary given fundamental domain, observe
that \(|\theta(z)|^2e^{-\pi H(z,z)}\) is periodic.
Its integral on any fundamental domain equals that on
the period parallelepiped.  Pull back by the real period
map; the Jacobian divided by
\(\operatorname{vol}\mathcal F\) is \(1\).
For \(z=x+\tau y\), the Gaussian in
\eqref{eq:ftb-key-gaussian-gauge} cancels exactly against
the Hermitian weight:
\[
 |e^{\pi z^tY^{-1}z/2}|^2e^{-\pi z^tY^{-1}\bar z}
                                     =e^{-2\pi y^tYy}.
\]
Lemma~\ref{lem:ftb-key-canonical-theta-integral} now yields
\[
 n_\theta=\frac{|c|^2}{\sqrt{2^g\det Y}},\qquad
 |\det W|=(2\pi)^{-g}|\det M|\,|\det\phi|,
 \qquad
 n_\eta=\frac{2^{2g}\pi^g\det Y}{|\det M|^2}.
\]
Substitution cancels \(c,M,Y,\phi\) and gives
\[
 \frac{|\det M|^4}{|c|^8}
 \frac{|c|^8}{2^{2g}(\det Y)^2}
 \frac{2^{4g}\pi^{2g}(\det Y)^2}{|\det M|^4}
                                      =(2\pi)^{2g}.
\]
This proves \eqref{eq:ftb-key-archnorm}.
\end{proof}

\paragraph{Source concordance.}
The prefix is
\srcid{src/Development/AbelianVarietyTheory/Proofs/}.
The exact Parseval--Tonelli--Gaussian calculation is
\srcid{ThetaLawProofs/LintegralLintegralNormSqRiemannTheta.lean}
(\srcid{bc5c74e3da3678aebab698a6e0c910788413b4aa}).
The coordinate, frame and constant calculation is
\srcid{ThetaLawProofs/NormalFormOfParts.lean}
(\srcid{7a587b5d72b8e340759219538bd7309c059f67de}).
Its real Jacobian and arbitrary-domain integral are
\srcid{VolumeFundamentalDomainPeriodLattice.lean}
(\srcid{fc3df3e87f65792bcfcd68064ad91c13e4dd9232})
and \srcid{SetIntegralPeriodLatticeEqVolumeMulCube.lean}
(\srcid{d0ba53fd9faf1c4a297671e128da7eafd83e0af4}),
in the same directory.
The sign-semicharacter calculation is
\srcid{CubicalLineBundles/SemicharacterSqEqOneOfSymmetric.lean}
(\srcid{d2d717aec9d26a9ed8ca968455911d6367e362bf}).
The actual section reading and half-characteristic recurrence
are \srcid{ThetaLawProofs/IsThetaFnAndNeZeroOfIsFn.lean}
(\srcid{462438291a306afec2495ae8147069869bdbc63b})
and \srcid{ExistsEqMulCanonicalThetaOfSignMultiplier.lean}
(\srcid{7d1ed10a5898cb126a40b8b262e087f0f7bd0bef}).
The actual arbitrary-uniformization comparison retained in
the first paragraphs is
\srcid{UnimodularAndTangentOfIsAH.lean}
(\srcid{a4df7948258a8441bcf1d2ffcd6db4a65fb44846}).
The final exact wrapper is
\srcid{KeyInNormalFormAndArchimedeanConstant.lean}
(\srcid{478dd63ff81fbea5390a62869ca1dc4e69706e91}).
The fixed-fibre analytic realization and its actual
linear-lift/Appell--Humbert comparison must be joined with
this calculation; the norm formula does not create those
geometric inputs.
